\chapter{Rationality in Economics: Axioms, Auxiliary Assumptions and Evidence}
\label{sec:p19-rationality-9-4}

Dhami and al-Nowaihi's 2021 chapter asks which meanings of economic
rationality are formal assumptions, which additional hypotheses are used
to obtain predictions, and which are supported by behavior. Its sequence
is important: specify preferences under certainty, risk, dated outcomes
and strategic interaction; separate auxiliary assumptions; then confront
both kinds of assumption with evidence. Precision and tractability make
a theory testable, but do not establish descriptive adequacy
\cite[pp.~565--574]{DhamiAlNowaihi2021Economics}.

This treatment preserves that comparison without identifying consistency
of preference with logical validity, empirical frequency with a normative
obligation, or a numerical utility representation with a cognitive process.
The chapter is a selective review, not a common experimental dataset.
Its examples have different tasks and evidential standing. A bounded
original jam-display study below makes those differences concrete.

\begin{hbmethod}{def:p19-r94-target}{Choose the target of a theory--evidence comparison}
Specify the alternatives, the comparison or choice task, the preference
relation and its domain, the information supplied, the model prediction,
and the observation that would contradict it. Keep these targets separate:
\begin{enumerate}
\item Axioms constrain a relation: completeness, transitivity, mixture
      independence or stationarity.
\item A representation theorem derives a numerical form under stated
      axioms. It does not establish that a participant satisfies them.
\item Auxiliary assumptions specify relevant consequences, beliefs,
      attention, computation, learning, institutions or market conditions.
\item An empirical record specifies a population, recruitment and task,
      design, response measures, uncertainty and generalization limits.
\end{enumerate}
An observed departure from a joint prediction can reject the conjunction
of its premises without identifying which premise failed. A normative
recommendation needs an independently defended goal and consequences.
Use \cref{def:p19-evidence-record} for the empirical record.
\end{hbmethod}

\section{Certainty: a relation and its representation}
\label{sec:p19-r94-certainty}

\begin{hbdefinition}{def:p19-r94-certainty}{The chapter's basic rationality criterion}
Let \(X\) be a nonempty set of alternatives. A weak preference
\(\succeq\) is complete when every \(x,y\in X\) satisfy
\(x\succeq y\) or \(y\succeq x\), and transitive when
\(x\succeq y\succeq z\) implies \(x\succeq z\).
Define \(x\sim y\) by both weak comparisons and \(x\succ y\)
by \(x\succeq y\) but not \(y\succeq x\).
The chapter calls a preference satisfying completeness and transitivity
rational under certainty. Indifference is not an absence of comparison.

When \(X\subseteq\mathbb R^n\), use the relative Euclidean topology.
The chapter's continuity condition requires both
\(\{x:x\succeq y\}\) and \(\{x:y\succeq x\}\) to be closed
in \(X\) for every fixed \(y\). A utility representation is a function
\(u:X\to\mathbb R\) such that
\[
                 x\succeq y\quad\Longleftrightarrow\quad u(x)\geq u(y).
\]
These are definitions on a declared domain, not observations about a
particular consumer \cite[pp.~565--566]{DhamiAlNowaihi2021Economics}.
\end{hbdefinition}

\begin{hbtheorem}{thm:p19-r94-continuous-existence}{cited=DhamiAlNowaihi2021Economics}{Continuous ordinal representation on a Euclidean subspace}
On a nonempty \(X\subseteq\mathbb R^n\), a complete, transitive
preference with the stated closed-contour continuity has a continuous
real-valued utility representation. This is the chapter's Proposition~1.
Its existence assertion concerns the entire specified domain, not only
an observed finite list of choices.
\end{hbtheorem}

\begin{hbtheorem}{thm:p19-r94-continuous-converse}{proved=annex:p19.r94.continuous-converse}{The continuous-representation converse}
If a continuous \(u:X\to\mathbb R\) represents \(\succeq\),
then \(\succeq\) is complete, transitive and has closed weak contours.
Any strictly increasing transformation of \(u\) preserves its ordinal
comparisons. Utility representation alone does not ensure that a maximum
exists on an infinite feasible set: \(u(x)=x\) on \(X=(0,1)\)
is a continuous representation without a maximizing alternative.
\end{hbtheorem}

The converse reconstructs the chapter's Proposition~2. The complete
finite construction is at \cref{thm:p17-finite-representation}, with
ties and degenerate sets explicitly handled. That finite result does not
by itself prove the infinite-domain existence assertion above.
Ordinal utility differences do not measure strength of preference;
interpersonal comparisons require additional structure. Neither
completeness nor transitivity requires that \(X\) describe only the
individual's own consumption. A failure of a preference axiom is not
automatically a violation of a law of logic
\cite[pp.~565--566, 568]{DhamiAlNowaihi2021Economics}.

\section{Risk, uncertainty and the added lottery axioms}
\label{sec:p19-r94-risk}

The chapter fixes finite real outcomes \(x_1\leq\cdots\leq x_n\).
A simple lottery \(L=(x_1,p_1;\ldots;x_n,p_n)\) has
\(p_i\geq0\) and \(\sum_i p_i=1\). A compound lottery has
lotteries as branches. Certain outcomes are degenerate lotteries.
The relevant preference now compares lotteries, not just certain prizes
\cite[pp.~566--567, note~1]{DhamiAlNowaihi2021Economics}.

\begin{hbdefinition}{def:p19-r94-risk-axioms}{The source's risk assumptions with explicit reduction}
Besides order, the chapter requires a strictly better designated best
outcome than worst outcome, calibration of every lottery to some
probability mixture of those endpoints, and common-mixture independence.
For lotteries \(L_1,L_2,H\) and \(0<q\leq1\), the last condition is
\[
 L_2\succeq L_1\quad\Longleftrightarrow\quad
 qL_2+(1-q)H\succeq qL_1+(1-q)H.
\]
The endpoint \(q=0\) is excluded: both right-hand lotteries would be
\(H\), so the reverse implication would force every pair indifferent.
Mixture independence is not independence of random events.

Reduction identifies a finite compound lottery with its final prize law:
if branch \(j\) has probability \(a_j\) and its conditional prize
probability is \(p_{ji}\), the reduced probability of prize \(i\) is
\(\sum_j a_jp_{ji}\). Repeat this operation for each level of a finite
tree. Nonnegative weights summing to one ensure a probability law;
zero-weight branches contribute nothing. Requiring indifference between
that procedure and its reduced law is the preference assumption, not
just the probability identity. The chapter names reduction but leaves
its formal statement to note~2.
\end{hbdefinition}

Under the finite-prize expected-utility assumptions, the representation is
\[
                  EU(L)=\sum_{i=1}^n p_i u(x_i).
\]
The chapter's Proposition~3 states this representation and its converse.
The full finite proof, endpoint calibration, reduction and positive-affine
uniqueness are supplied at \cref{thm:p17-finite-vnm,thm:p17-lottery-reduction}.
The calibration requirement above is the source's continuity formulation;
the neighboring theorem states its precise mixture/topological assumptions.
The following bridge also derives the representation directly from the
chapter's endpoint-calibration formulation.

Strict mixture-segment ordering together with calibration makes the
designated endpoints best and worst for the entire lottery order; this
does not add a preference restriction to the source's calibrated domain.

\begin{hbtheorem}{thm:p19-r94-source-vnm}{proved=annex:p19.r94.source-vnm}{Representation from the chapter's calibrated finite lottery axioms}
For finite prizes, reduce every finite compound lottery to its final law.
Suppose the lottery order is complete and transitive, designated certain
lotteries \(M\succ m\) are its best and worst, every lottery is
indifferent to some \(qM+(1-q)m\) with \(q\in[0,1]\), and
common-mixture independence holds for positive mixture weights. There is
a unique normalized prize utility, equal to 1 and 0 on the respective
prizes underlying \(M\) and \(m\), such
that \(EU(L)=\sum_i p_i u(x_i)\) represents the order. All other
expected-utility prize representations are its positive-affine transforms.
Conversely, such an expected-utility representation with the stated
designated best and worst satisfies these axioms.
\end{hbtheorem}

Expected utility averages consequence utility, not monetary amounts unless
\(u\) is linear. Reduction omits procedure-dependent consequences unless
they have first been included in the prize description. Even an expected
utility model can value fairness or another person's consumption
\cite[pp.~566--568]{DhamiAlNowaihi2021Economics}.

In the chapter's uncertainty discussion, probabilities become subjective.
A preference representation under supplied risks does not establish either
the existence of a suitable subjective probability or its calibration.
The chapter assumes Kolmogorov probability and later defines Bayesian
rationality to add Bayesian updating. Those are separate assumptions,
not consequences of completeness and transitivity. The subjective-bet
and uncertainty treatments at
\cref{sec:p17-finite-subjective-probability,sec:p17-rationality-8-4}
retain their different domains and representation conditions
\cite[p.~567, note~4]{DhamiAlNowaihi2021Economics}.

\section{Dated outcomes and consumption streams}
\label{sec:p19-r94-time}

\begin{hbdefinition}{def:p19-r94-time-axioms}{Time preference on an admissible outcome--date domain}
Write \(Z\) for a nondegenerate real interval of outcomes and \(\Gamma\)
for either successive nonnegative integers or an interval of nonnegative
reals containing 0 and 1. A dated outcome is \((z,t)\in Z\times\Gamma\).
The five requirements are: complete and transitive ordering; strict
preference for larger outcomes at a fixed date; earlier receipt preferred
for positive outcomes, later receipt for negative outcomes, and date
indifference for zero if present; closed weak contours in the product
topology; and stationarity of indifference.

The last condition says that
\((z,t)\sim(w,t+\tau)\) holds exactly when
\((z,s)\sim(w,s+\tau)\), whenever every displayed date belongs to
\(\Gamma\). It applies to admissible common shifts, not to dates outside
the domain. These domain and loss qualifications are those of the original
Fishburn--Rubinstein A0--A5; the chapter's B1--B5 abbreviate them
\cite[pp.~680--681]{FishburnRubinstein1982TimeR94}
\cite[p.~567]{DhamiAlNowaihi2021Economics}.
\end{hbdefinition}

\begin{hbtheorem}{thm:p19-r94-time-representation}{cited=FishburnRubinstein1982TimeR94}{The original dated-outcome representation and its uniqueness boundary}
Under the stated outcome--date assumptions, for any fixed \(0<\delta<1\)
there is a continuous strictly increasing \(u:Z\to\mathbb R\) representing
the preference by
\[
       (z,t)\succeq(w,s)\quad\Longleftrightarrow\quad
                         \delta^t u(z)\geq\delta^s u(w).
\]
If zero is present, \(u(0)=0\); otherwise \(u(z)\) has the sign of
\(z\). For interval time, at fixed \(\delta\), uniqueness is up to
a positive multiplier on each nonempty positive and negative outcome
sector separately. The original theorem does not assert this uniqueness
for discrete time. This is Theorem~2 of Fishburn and Rubinstein, whose
published construction explicitly omits some details.
\end{hbtheorem}

Freedom to choose \(\delta\) in this single-outcome representation is
not identification of a psychological discount parameter: changing it
changes \(u\). Nor do axioms on one dated outcome alone prove the
additive stream formula \(\sum_t\delta^t u(c_t)\). That model also
assumes an aggregation across dates with appropriate separability.
Stationarity is a substantive descriptive restriction, not an implication
of ordinary preference order. The original authors themselves question
its psychological plausibility
\cite[pp.~681--686]{FishburnRubinstein1982TimeR94}.

\begin{hbdefinition}{def:p19-r94-stream}{The consumption-planning example and its feasibility conditions}
For a finite horizon \(T\), savings at date \(i\) evolve as
\[
        S_i=y_i+(1+r)S_{i-1}-c_i,\qquad r\geq0,
\]
where income \(y_i\), consumption \(c_i\), an initial stock, and
admissible consumption/borrowing sets must be supplied. With no bequest
utility and terminal solvency \(S_T\geq0\), a time-\(t\) objective is
\[
             \sum_{i=t}^T\delta^{\,i-t}u(c_i),\qquad 0<\delta<1.
\]
The positive forward exponent discounts later utility. The chapter's
printed reversed exponent would instead increase its weight.
Its instruction to consume all available final resources additionally
requires increasing utility and feasibility of that consumption; solvency
alone does not imply a tight terminal constraint.
\end{hbdefinition}

\begin{hbtheorem}{thm:p19-r94-backward}{proved=annex:p19.r94.backward}{Backward induction for the finite planning version}
Suppose the reachable states and action sets are finite and every retained
state has a feasible continuation. Define terminal continuation value to
be zero at solvent terminal states and remove insolvent continuations.
At date \(i\), let
\[
 V_i(s)=\max_{c\text{ feasible at }(i,s)}
        \{u(c)+\delta V_{i+1}(y_i+(1+r)s-c)\}.
\]
This recursion gives the largest feasible discounted utility and a policy
attaining it. It proves optimization in the supplied finite model, not
that people implement the algorithm or that every continuous model has
an attained solution.
\end{hbtheorem}

\begin{hbexample}{def:p19-r94-planning-example}{An explicit finite planning contrast}
Take two dates, \(S_0=2\), no income, \(r=0\), no borrowing,
\(u(c)=\sqrt c\), \(\delta=1/2\), and integer consumption. The
feasible terminal-exhausting allocations \((0,2),(1,1),(2,0)\)
give \(\sqrt2/2,3/2,\sqrt2\), respectively. The middle allocation
is best among this finite set. This stipulated example illustrates the
conditional recursion, not a measured household preference or advice
about saving.
\end{hbexample}

The source's forty-year example points to \(2^{40}>10^{12}\) binary
income histories, contrasting a three-year enumeration with eight
histories. This identifies growth of a history list, not a necessary
number of operations for every solution method: repeated states,
structure, approximation or closed forms can change complexity.
The chapter reports that simplified experiments need not show backward
induction and emphasizes uncertain lifespan, income and returns,
bequests, emotional states and imperfect capital markets. Its survey
does not supply a single task/population estimate for that claim
\cite[pp.~568--569]{DhamiAlNowaihi2021Economics}.

\begin{hbdefinition}{def:p19-r94-present-bias}{A different stream model}
For positive \(\delta<1\) and \(0<\beta<1\), quasi-hyperbolic
evaluation at the present is
\[
 U(c_0,\ldots,c_T)=u(c_0)+\beta\sum_{t=1}^T\delta^t u(c_t).
\]
The extra \(\beta\) discounts all future utility relative to present
utility. At a later date, a new current outcome loses that extra factor.
Awareness of future preferences and the availability of commitment are
further modeling choices, not determined by this formula.
\end{hbdefinition}

\begin{hbtheorem}{thm:p19-r94-time-contrast}{proved=annex:p19.r94.time-contrast}{A stationarity contrast with declared values}
An exponential dated-outcome comparison is unchanged when both dates
receive the same admissible shift. With linear \(u\), \(\delta=.9\)
and \(\beta=.5\), quasi-hyperbolic evaluation prefers 10 now to
15 one date later, but prefers 15 at date 2 to 10 at date 1 when both
are viewed from date 0. Re-evaluation at date 1 reverses that latter
ranking. These are model calculations, not measured frequencies.
\end{hbtheorem}

The source reports Thaler's monetary-indifference examples: 15 now versus
60 after a year, 250 versus 350, and 3000 versus 4000; and 15 now versus
20 after a month, 50 after a year and 100 after ten years. Under a
linear monetary-utility interpretation, annual factors are respectively
\(1/4,5/7,3/4\), and \((3/4)^{12},3/10,(3/20)^{1/10}\).
They are approximately .25, .71, .75 and .032, .30, .827.
The chapter also reports smaller discounting of losses than gains and
preferences for shapes of consumption profiles
\cite[p.~571]{DhamiAlNowaihi2021Economics}.

These are the chapter's review-reported illustrations, not an independently
reconstructed Thaler dataset. It does not give the participants, full
elicitation procedure, distribution or interval uncertainty for these
numbers. The ratios do not identify discount factors without assumptions
about utility, receipt risk, transaction costs and timing of consumption.
Magnitude effects cannot be inferred from money ratios alone when \(u\)
is nonlinear. Likewise, claims about consumption tracking income,
retirement consumption drops, simultaneous illiquid assets and debt,
addiction, cigarette-pack size or taxes, and gym contracts are proposed
applications of the alternative model, not consequences proved by those
few examples \cite[p.~571, note~7]{DhamiAlNowaihi2021Economics}.

\section{Auxiliary assumptions are not preference axioms}
\label{sec:p19-r94-auxiliary}

The chapter separates four influential auxiliary assumptions. First,
self-regarding preferences value own consumption only; social preferences
can value the allocation \((x_i,x_{-i})\). Completeness and transitivity
can hold on either domain. Dictator, ultimatum, public-goods, trust and
gift-exchange games test different social-preference mechanisms; their
names do not define one shared measure of altruism. The review attributes
its social-preference evidence to the surveyed literature, including
laboratory, field and neuroeconomic records
\cite[p.~568, note~3]{DhamiAlNowaihi2021Economics}.

The authors contrast the centrality of social preferences in behavioral
economics with their historically smaller role in mainstream modeling,
while noting exceptions such as keeping-up-with-the-Joneses and snob
consumption. This is their assessment of modeling practice, not a
restriction imposed by the representation theorem.

\begin{hbexample}{def:p19-r94-social-emotional}{Consistent preferences need not be selfish or emotionless}
On two-person nonnegative allocations, let \(u(x_1,x_2)=x_1+x_2\).
The induced weak order prefers \((1,3)\) to \((2,0)\) for person 1,
although own consumption is lower. It is complete and transitive.
Alternatively let an emotional state \(m\) select one of two complete
orders, each represented by its own \(u_m\). Different rankings across
states need not be a cycle within either state's order. These are
counterexamples to implications, not evidence that a particular person
has either function.
\end{hbexample}

Second, the absence of emotions from a baseline representation is a
restriction on what is modeled, not a theorem that rational choice cannot
involve emotion. The chapter proposes indexing preferences by emotional
state and discusses anger as a credible signal, anxiety in anticipatory
utility, projection of current feelings onto future circumstances, and
guilt or intention sensitivity in psychological game theory. Its kidney-
transplant/well-being example is a reviewed comparison, without a primary
sample, outcome scale or interval record in this chapter; it cannot alone
identify projection bias or establish a causal effect of transplantation
\cite[p.~568]{DhamiAlNowaihi2021Economics}.

Third, the source criticizes perfect attention, unlimited calculation and
absence of misperception. An abstract preference relation need not have
been produced by explicitly calculating every pairwise comparison.
An optimization prediction also need not assert that a consumer consciously
runs the optimizer; an ``as if'' account and a literal process account
require different tests. Yet a process hypothesis invoking substantial
computation must explain resources, timing and feasibility. The source's
planning example exposes this issue, especially under distraction,
deadlines and information overload
\cite[pp.~568--569]{DhamiAlNowaihi2021Economics}.

Fourth, cancellation of errors and elimination of noise traders are additional
market hypotheses. If every participant's signed valuation error is +1,
their average error is also +1; averaging does not force cancellation.
Whether arbitrage corrects an error depends on beliefs, funding, risks,
constraints and horizon. This analytical counterexample does not estimate
any market's actual error \cite[pp.~569--570, 573--574]{DhamiAlNowaihi2021Economics}.

\section{A bounded original choice record}
\label{sec:p19-r94-jam}

\begin{hbdefinition}{def:p19-r94-jam-record}{Iyengar--Lepper Study 1: display, attraction and purchase}
The 2000 field study used a tasting booth in Draeger's, an upscale Menlo
Park grocery store, on two consecutive nonholiday Saturdays, each for
five hours. Six versus twenty-four exotic Wilkin \& Sons jam flavors
were displayed in hourly rotations, counterbalanced across days. A
thirty-student preliminary preference survey selected the six-flavor
subset to include highly, moderately and poorly rated flavors. This was
not individual random assignment of shoppers to conditions.

Approximately 754 shoppers were observed: 386 were present during the
extensive-display hours and 368 during limited-display hours. Actual
encounters numbered 242 and 260; booth stops numbered 145 and 104.
Those who stopped were observationally classified as mostly middle-aged
and Caucasian, with about 62 percent women. These classifications are
not a probability-sampled population frame. Each visitor received a
one-dollar coupon, valid for a week. Anyone purchasing then encountered
the entire shelf assortment, regardless of the earlier booth condition.

The measures were booth attraction, number of tastes and coupon-linked
purchase, not pairwise preference completeness, subjective utility,
deliberation time or an inference correctness score. A blinded rating
check on recorded solicitations assessed friendliness, rather than
measuring the hypothesized decision mechanism
\cite[pp.~996--997]{IyengarLepper2000ChoiceR94}.
\end{hbdefinition}

Sixty percent of extensive-display encounterers stopped (145 of 242),
versus forty percent of limited-display encounterers (104 of 260);
the reported attraction comparison was
\(\chi^2(1,N=502)=19.89,p<.001\). Among those stopping, purchase
was 31 of 104 in the limited condition versus 4 of 145 in the extensive
condition, approximately 29.8 versus 2.8 percent; the reported comparison
was \(\chi^2(1,N=249)=32.34,p<.0001\). Attraction and purchase
therefore have different denominators. Mean numbers sampled were 1.38
and 1.50, respectively; the paper reported no significant difference,
\(F(1,245)=.83\). This concerns counts of tastes per visitor, not equal
percentages of all available flavors sampled
\cite[p.~997]{IyengarLepper2000ChoiceR94}.

The printed comparisons provide significance tests, not confidence intervals
for the purchase contrast. Visitors self-selected into stopping; the larger
display might attract otherwise unmotivated shoppers. The authors also
consider perceived special selection of the six jams and social/time
constraints on tasting. Hourly counterbalancing does not create independent
individual randomization, and the selected-visitor purchase contrast is
not the same as a population causal effect of menu size. The store,
product, coupon and unfamiliar-flavor selection constrain transfer
\cite[pp.~997--998, 1003--1004]{IyengarLepper2000ChoiceR94}.

The selected methods do not report exact collection calendar dates,
repeat-visitor identification or a full missing-observation audit. They
also do not supply a manipulation measuring the proposed completeness
mechanism independently of purchase.

The printed purchase counts also leave a numerical reporting discrepancy.
Their two-by-two table is \((31,73;4,141)\), where each row gives
purchase and nonpurchase among visitors in its condition. Ordinary Pearson
calculation on that table gives approximately 36.68, rather than the
published 32.34. The attraction table reproduces its reported 19.89.
The published purchase test is retained as reported; its mismatch is not
resolved by the chapter or this calculation, and no raw-data correction
is inferred. The design limitations remain even if either test is highly
significant.

The 2021 chapter interprets this as possible difficulty choosing with more
options and a possible challenge to completeness. The difficulty interpretation
is distinct from an identified violation: declining to buy, searching longer
or accepting an outside option can occur under a complete preference,
especially when information and decision costs differ. The original authors
explicitly leave mechanism and moderation questions open. Their discussion
does not support a universal six-option rule, nor a welfare policy based
solely on sales \cite[p.~570]{DhamiAlNowaihi2021Economics}
\cite[pp.~998, 1003--1004]{IyengarLepper2000ChoiceR94}.

\section{Reversals and risk: what is actually contradicted?}
\label{sec:p19-r94-reversals}

The source's preference-reversal task compares
\(P\text{-bet}=(0,1-P;z,P)\) and
\(\$\text{-bet}=(0,1-p;Z,p)\), where \(p<P\) and \(z<Z\).
One alternative is more likely to win; the other has a larger prize.
Direct choice and certainty-equivalent elicitation are distinct response
tasks. A certainty equivalent is a certain consequence indifferent to a
lottery, not merely a stated selling price under unspecified incentives
\cite[p.~570, note~5]{DhamiAlNowaihi2021Economics}.

\begin{hbtheorem}{thm:p19-r94-reversal}{proved=annex:p19.r94.reversal}{A reversal contradicts a single transitive order only under a common interface}
Suppose one relation compares both bets and certain monetary consequences,
\(P\text{-bet}\succ\$\text{-bet}\),
\(P\text{-bet}\sim C_P\), \(\$\text{-bet}\sim C_{\$}\), and
\(C_{\$}\succ C_P\). These comparisons cannot belong to one
complete transitive preference. Interpreting \(C_P<C_{\$}\) as that
strict comparison additionally assumes monotonicity for money with all
other consequences held fixed.
\end{hbtheorem}

The chapter reports the pattern as replicated in the preference-construction
literature, but supplies neither a shared population nor effect interval.
The formal contradiction is conditional on the common relation and genuine
indifference. Distinct frames, elicitation noise, incentives or consequences
can instead make the records comparisons of different objects; that is an
alternative explanation to test, not permission to redefine every failure
away. The source itself notes that narrow frames may preserve local
consistency without explaining all contrary evidence
\cite[p.~570]{DhamiAlNowaihi2021Economics}.

\begin{hbexample}{def:p19-r94-allais}{A declared common-consequence contrast}
In arbitrary prize units, let \(A\) give 1 certainly, \(B\) give
1 with probability .89, 5 with probability .10 and 0 with probability .01,
\(C\) give 1 with probability .11 and 0 otherwise, and \(D\) give
5 with probability .10 and 0 otherwise. The pattern \(A\succ B\)
and \(D\succ C\) is a classic Allais-type contrast. These particular
numbers specify a worked mathematical illustration, not a new experiment
or an estimated frequency.
\end{hbexample}

\begin{hbtheorem}{thm:p19-r94-allais}{proved=annex:p19.r94.allais}{The two strict comparisons cannot both have one expected-utility representation}
No single prize utility \(u\), evaluated linearly in the given
probabilities, can produce both strict comparisons in the illustration.
No assumption of linear utility in monetary prizes is needed.
\end{hbtheorem}

The chapter identifies Allais-type independence violations, nonlinear
probability weighting, rank-dependent utility and prospect theory as central
to the descriptive comparison. Nonlinear probability weighting is a change
in evaluation, not simply an alternative probability measure. Its broader
review families include risk attitudes, endowment effects, tax evasion,
goal setting, taxi labor supply, equity premia, contracting and asset returns.
Those observations do not all test independence under the same design;
reference points, consequence descriptions and auxiliary assumptions matter.
The prospect-theory treatment at \cref{sec:p17-rationality-8-3} preserves
the distinctions between versions, weighting, curvature and reference
dependence \cite[pp.~570--571, note~6]{DhamiAlNowaihi2021Economics}.

The source also invokes capuchin behavior as resembling prospect-theory
patterns and interprets it as inherited from a common ancestor. Behavioral
resemblance alone does not establish genetic inheritance of a particular
representation, exclude learning, or identify identical task mechanisms
across species. No cross-species sample, task equivalence or inheritance
design is supplied in this chapter; the interpretation remains a stronger
claim than the reported resemblance \cite[p.~571]{DhamiAlNowaihi2021Economics}.

\section{Strategic predictions and their evidential interfaces}
\label{sec:p19-r94-games}

A Nash equilibrium specifies mutual best responses under the game's
payoffs and distributions. Its mathematical definition does not assert
that players' actual beliefs satisfy particular higher-order knowledge
conditions. The source discusses demanding epistemic conditions for an
equilibrium conjecture and distinguishes this justification from learning.
The actual two-player and many-player Aumann--Brandenburger conditions,
common-prior/null-type boundaries and conjecture-versus-action distinction
are at \cref{sec:p17-epistemic-nash-conditions}; evolutionary stability and
learning dynamics are different questions at
\cref{sec:p17-rat93-evolutionary-games}
\cite[pp.~567--568]{DhamiAlNowaihi2021Economics}.

Failure of an empirical equilibrium prediction does not isolate the failure
of its epistemic conditions unless the observed action, conjecture and
experimental game genuinely satisfy the theorem's interfaces. The chapter
also warns that a learning process need not converge, plausible convergence
need not be rapid, and many important choices provide little repeated
experience. Early-round and unfamiliar-game behavior are therefore part of
the intended descriptive target, not automatically expendable warm-up data
\cite[pp.~568, 571--572]{DhamiAlNowaihi2021Economics}.

The chapter's ten strategic-evidence families should be distinguished:
\begin{enumerate}
\item Limited depths of iterated dominance, rather than arbitrary failure
      of all strategic reasoning.
\item Different responses to extensive and corresponding normal forms,
      rather than formal inequivalence of the represented game.
\item Backward-induction departures in dynamic games, including the review's
      centipede/chess-expert illustration.
\item Difficulty implementing even selected dominance-solvable mechanisms,
      without an inference that every more complex mechanism must fail.
\item Mixed play not matching equilibrium proportions; the review retains
      some supportive sports findings while questioning confounds and transfer.
\item History-sensitive equilibrium selection and improved coordination
      with preplay communication, not one universal selection principle.
\item Bargaining counteroffers, impasse and asymmetric-information contrasts;
      the source retains unexplained support in its two-sided-information
      comparison instead of suppressing the favorable case.
\item MOUSELAB information-search patterns in alternating-offers bargaining,
      which concern process as well as terminal choices.
\item Signaling predictions and the reviewed boundary around the
      Cho--Kreps intuitive criterion.
\item Guilt, surprise, anger and inferred intentions as possible payoff
      ingredients in psychological games.
\end{enumerate}
These are the 2021 authors' survey claims. The chapter sends the reader
to Camerer and Dhami for the underlying experiments and does not supply
their individual populations, treatments, measures or uncertainty. No one
measured effect or general rejection rate is inferred from this list
\cite[pp.~571--572]{DhamiAlNowaihi2021Economics}.

Level-\(k\), quantal-response, analogy-based and cursed-equilibrium accounts,
team reasoning and evidential reasoning are named alternatives in the
source. They relax different assumptions and are not interchangeable labels
for ``irrationality.'' Correctly diagnosing a mechanism requires the specific
game, beliefs, incentives and process observations.

\section{Attention, heuristics and bounded methods}
\label{sec:p19-r94-bounded}

The source's limited-attention examples concern tax salience at posted
prices versus checkout, electronic versus manual tolls, and used-car mileage
near the 9999/10000 boundary. It also discusses scarcity-related attention
demands and formal sparse-attention models. These are distinct environments
and measures; the chapter supplies no common causal effect or sample
covering all of them. An omission from a consideration set is not by itself
an incompleteness observation on an otherwise fixed preference domain
\cite[p.~570]{DhamiAlNowaihi2021Economics}.

The heuristics discussion contrasts full Bayesian optimization with regular,
nonrandom departures from it. Representativeness can produce excessive
expectation of small samples resembling their population; proposed random
sequences can show too many alternations. Gambler's-fallacy and hot-hand
interpretations require a specified independent-outcome benchmark: detecting
dependence in an actually dependent process is not that fallacy. The source
also names availability, anchoring, base-rate neglect, conservatism,
hindsight, confirmation and regression-related misjudgments. These labels
do not establish a shared process or a universal direction of error
\cite[pp.~572--573, notes~8--9]{DhamiAlNowaihi2021Economics}.

The chapter acknowledges favorable heuristic performance under uncertainty
about an optimal solution, while criticizing unclear optimization benchmarks.
That qualification matters. \Cref{sec:p17-rationality-8-5} distinguishes
resource-constrained optimization, satisficing and ecological procedures;
\cref{sec:p19-rationality-2-4} separates outcome bias from process and
environment. Less information need not be worse, but a heuristic's success
requires its actual information environment and loss. It does not imply
that no model-based method can improve on it.

Frequency representation is treated cautiously in the source: it can help
some tasks without eliminating every heuristic-related error. The natural-
frequency experiments at
\cref{sec:p19-natural-frequencies-bayesian-reasoning} preserve nested-set
structure, event interpretation and response format; they cannot be
transferred merely by replacing percentages with counts. The chapter's
claim about what format most real economic information uses is not a
measured prevalence estimate in its evidence record
\cite[p.~573]{DhamiAlNowaihi2021Economics}.

Coarse categories and mental models offer other explanations of how people
reduce cognitive demands. The source connects them to analogy-based and
cursed equilibria, persuasion, causal beliefs, social identities and cultural
worldviews; models may help or hinder depending on accuracy and persistence.
These uses of ``mental model'' are not automatically the same experimental
construct as the inference models at \cref{sec:p19-rationality-2-3}.
The chapter's mixture of developmental, cultural and evolutionary claims
does not establish one origin for all such representations
\cite[p.~573]{DhamiAlNowaihi2021Economics}.

Simon-style satisficing searches sequentially and stops when an aspiration
criterion is met; it can explain behavior without claiming exhaustive
maximization. The source cites a search experiment and an aspiration-based
cooperation model as support, while distinguishing fast-and-frugal work
that often uses experimenter-supplied heuristics. A formal cooperation
construction is not the same as an observed population effect. Its
computational warning also needs the proper complexity distinction:
NP-completeness is a worst-case classification of decision problems, not
a proof that every economic problem or every algorithm requires exponential
time. The traveling-salesman example names the route optimization; its
decision version asks whether a tour under a given length bound exists
\cite[p.~573, note~10]{DhamiAlNowaihi2021Economics}.

\section{Markets, generalization and the limits of the review}
\label{sec:p19-r94-markets}

The chapter's efficient-market discussion combines Bayesian updating,
fundamental-value pricing, error cancellation and arbitrage survival.
Its fundamental value is a discounted future-income valuation, which
requires an information set, cash-flow expectations, discounting and a
return/risk model before a deviation is measurable. Information efficiency,
one particular fundamental-value model and the behavior of each trader
are different hypotheses. A pricing anomaly generally tests a joint
specification rather than completeness and transitivity alone
\cite[pp.~569--570, 573--574]{DhamiAlNowaihi2021Economics}.

The reviewed families are persistent price differences between substitute
assets, gradual information incorporation, short-horizon momentum and
long-horizon reversal, bubbles, crashes described as without news, and
under- or overreaction. The source offers correlated noise-trader errors,
beliefs that mispricing will grow, and performance-linked withdrawals
forcing premature liquidation as mechanisms limiting corrective arbitrage.
Institutional participation therefore does not logically guarantee
correction. Its CEO-overconfidence/merger illustration is another distinct
corporate-finance claim, not an estimated merger-failure effect supplied
by this chapter \cite[pp.~573--574, note~11]{DhamiAlNowaihi2021Economics}.

The source also attributes pre-crisis nonintervention to Greenspan's belief
in market efficiency and describes consequences as disastrous. This is
the authors' historical and causal interpretation, not an identified
policy counterfactual or a primary event record in this chapter. Neither
that interpretation nor a catalog of anomalies supplies an investment
strategy or predicts a particular asset's future price.

\begin{hbmethod}{def:p19-r94-use}{Use the comparison without overgeneralizing it}
For a preference-axiom claim, test the same alternatives and relation
with controlled interpretation, including ties and outside options. For
an expected-utility or discounting claim, distinguish prize utility from
probability/date weights and state the elicitation and identification
assumptions. For a strategic prediction, distinguish actions from beliefs,
equilibrium from convergence, and early-round from trained behavior.
For a market claim, state the information and return model and constraints
on arbitrage. Retain recruitment, task, design, measures and uncertainty
before transferring any empirical result to a new population or setting.

A successful use is the conditional reversal contradiction under one
declared relation, or the original jam record with its two denominators
and selection limitations. A failed use is treating lower jam purchases
as proof of incomplete preferences, assigning an annual discount factor
from money ratios without a utility model, or using a theorem's hypothetical
belief conditions as measured facts about an experimental participant.
\end{hbmethod}

\begin{proofboundary}{def:p19-r94-boundary}{Formal and empirical standing}
The chapter concludes that neoclassical precision is valuable but does
not shield its rationality and auxiliary assumptions from contrary
evidence, and that behavioral economics develops rigorous alternatives.
This is its argument, not a single universal empirical rejection result
\cite[p.~574]{DhamiAlNowaihi2021Economics}.

The general continuous-utility existence assertion and original dated-
outcome representation retain their actual cited boundaries. Existing
finite representation and epistemic proofs retain their own hypotheses;
the six new reciprocal annexes prove the converse, source-calibrated
expected-utility representation, finite recursion,
time contrast and two conditional risk/reversal contradictions. The
original jam record is bounded to its actual task and sample. Other
literature families remain attributed review reports whose study-level
design/uncertainty fields are not supplied by the chapter; they cannot
discharge a separately required original-study or full-proof obligation.
No universal preference, process, welfare, policy or investment prescription
is derived from their aggregation.
\end{proofboundary}
